\section{Main concepts}

In this section, we present the generalized maps, graph transformation rules and their subsequent concepts which are necessary to the understanding of our contribution.

\subsection{Generalized maps}\label{ssec:g-maps}

\emph{Generalized maps} (or G-maps) \cite{lienhardt1991topological, Damiand-Lienhardt14} allow the representation of manifold geometric objects (with or without boundaries), based on some cellular $n$-dimensional topological structure.

\begin{figure*}[t]
    \begin{subfigure}{.19\textwidth}
        \centering
        \includegraphics[height=.18\textheight]{decomposition-3D-no-decomp}
        \caption{}
        \label{fig:demo-fig-ABCDE}
    \end{subfigure}
    \hfill
    \begin{subfigure}{.19\textwidth}
        \centering
        \includegraphics[height=.18\textheight]{decomposition-3D-3-cell}
        \caption{}
        \label{fig:decomp-volume}
    \end{subfigure}
    \hfill
    \begin{subfigure}{.19\textwidth}
        \centering
        \includegraphics[height=.18\textheight]{decomposition-3D-2-cell}
        \caption{}
        \label{fig:demo-fig-faces-decomposition}
        \end{subfigure}
    \hfill
    \begin{subfigure}{.19\textwidth}
        \centering
        \includegraphics[height=.18\textheight]{decomposition-3D-1-cell}
        \caption{}
        \label{fig:demo-fig-show-1-links}
    \end{subfigure}
    \hfill
    \begin{subfigure}{.19\textwidth}
        \centering
        \includegraphics[height=.18\textheight]{decomposition-3D-0-cell}
        \caption{}
        \label{fig:demo-fig-show-all-links}
    \end{subfigure}
    \caption{Cell decomposition of a geometric 3D object}
    \label{fig:demo-fig-cell-decomposition}
\end{figure*}

\maxime{Refaire le paragraphe à partir des nouvelles figures}

The representation of an object as a G-map comes intuitively from its decomposition into topological cells (vertices, edges, faces, volumes, and so on).
For example, the 2D topological object (Fig.~\ref{fig:house-ABCDE}) can be decomposed into a $2$-dimensional G-map (or $2$-G-map).
The object is first decomposed into faces (Fig.~\ref{fig:house-faces-decomposition}).
These faces are \emph{linked} along their common edge with an $2$ link, drawn in blue.
Each edge on the border of the object is connected to itself (the blue $2$ link forms a loop).
The index "$2$" means that the link connect two $2$-dimensional (possibly a single one) faces.
In order to simplify the reading, $2$ labels will not be written
in every subsequent figure.

In the same way, faces are split into edges connected with red $1$ links (Fig~\ref{fig:house-show-1-links}).
Lastly, edges themselves are split into vertices with black $0$ links (Fig~\ref{fig:house-show-all-links}) to produce the $2$-G-map describing the objects shown in Fig.~\ref{fig:house-ABCDE}.
A G-map is therefore a graph, the nodes of which are the vertices (named \emph{darts}) obtained at the end of the process and the edges are $i$ links.


G-maps have conditions guaranteeing objects consistency, %coherence,
for example, two faces are always linked along an edge.

\begin{figure*}[t]
  \begin{subfigure}{.19\textwidth}
    \centering
    \includegraphics[height=.18\textheight]{orbit-3D-vertex-123.pdf}
    \caption{}
    %\label{fig:house-2-cell-a}
    \label{fig:orbit-3D-vertex-123}
  \end{subfigure}
  \hfill
  \begin{subfigure}{.19\textwidth}
    \centering
    \includegraphics[height=.18\textheight]{orbit-3D-edge-023.pdf}
    \caption{}
    %\label{fig:house-1-cell-a}
    \label{fig:orbit-3D-edge-023}
  \end{subfigure}
  \hfill
  \begin{subfigure}{.19\textwidth}
    \centering
    \includegraphics[height=.18\textheight]{orbit-3D-face-013}
    \caption{}
    %\label{fig:house-1-cell-c}
    %\label{fig:orbit-3D-face-013}
  \end{subfigure}
  \hfill
  \begin{subfigure}{.19\textwidth}
    \centering
    \includegraphics[height=.18\textheight]{orbit-3D-volume-012}
    \caption{}
    %\label{fig:house-0-cell-a}
    \label{fig:orbit-3D-volume-012}
  \end{subfigure}
  \hfill
  \begin{subfigure}{.19\textwidth}
    \centering
    \includegraphics[height=.18\textheight]{orbit-3D-edge-0}
    \caption{}
    %\label{fig:house-1-cell-a0-a}
    \label{fig:orbit-3D-edge-0}
  \end{subfigure}

  \caption{Orbit decomposition of a geometric 3D object \maxime{TODO: réorienter + transparence tt le temps}}
  \label{fig:house-orbit-decomposition}
\end{figure*}



Topological cells are not explicitly represented in G-maps but only implicitly defined as subgraphs. They can be computed using graph traversals defined by an originating node and by a given set of link labels named \emph{orbit}.
For example, the 2-cell adjacent to some dart $a$ (or the object face matching $a$) (Fig.~\ref{fig:house-2-cell-a}) is the subgraph which contains $a$ and all darts reachable from $a$, using links labeled by $0$ or $1$ (\textit{i.e.} darts $a$, $b$, $c$, $d$, $e$, $f$, $g$ and $h$) and the links themselves.
This subgraph is denoted by $G\orbit{0, 1}(a)$ and models the face $BCED$ (Fig.~\ref{fig:house-ABCDE}).
$\orbit{0, 1}$ is the \emph{type} of the orbit.
The $1$-cell adjacent to $a$ (or the object edge matching to  $a$) (Fig.~\ref{fig:house-1-cell-a}) is the subgraph $G \orbit{0, 2}(a)$ which contains $a$ and all reachable darts using links labeled by $0$ or $2$ (\textit{i.e.} darts $a, b, l$ and $m$) and the corresponding links.
It represents the topological edge $BC$.
Thanks to $2$-loops on the object border, $G\orbit{0, 2}(c)$ (Fig.~\ref{fig:house-1-cell-c}) also represents the edge $CE$ matching dart $c$.
The $0$-cell adjacent to $a$ (or the object vertex matching dart $a$) (Fig.~\ref{fig:house-0-cell-a}) is the subgraph $G\orbit{1, 2}(a)$ and represents vertex $B$.
Note that orbits are more general than cells.
For example, the face edge $G\orbit{0}(a)$ (Fig.~\ref{fig:house-1-cell-a0-a}) is an $\orbit{0}$ orbit incident to $a$.


\subsection{Graph transformation rules}
\label{ssec:rules}

Jerboas's \cite{belhaouari2014jerboa} graph transformation rules allow the formalisation of operations over G-maps.

\begin{figure*}[t]
  \centering
  % \maxsizebox{\columnwidth}{\columnwidth}{}
  \includegraphics[width=0.9\textwidth]{rule-face-extrusion} \caption{Extrusion rule}
  % \label{fig:rule-face-triangulation}
  \label{fig:rule-face-extrusion}
\end{figure*}

\begin{figure}[t]
  \centering
  \begin{subfigure}[t]{.45\columnwidth}
    \centering
    \includegraphics[width=.8\textwidth]{demo-fig-extrusion-left-match.pdf}
    \caption{Graph $G$}
    \label{fig:demo-fig-triangulation-left-match}
  \end{subfigure}
  \hfill
  \begin{subfigure}[t]{.45\columnwidth}
    \centering
    \includegraphics[width=.8\textwidth]{demo-fig-extrusion-right-match.pdf}
    \caption{Graph $H$}
    \label{fig:demo-fig-triangulation-right-match}
  \end{subfigure}
  \caption{Application of the extrusion rule (Fig.~\ref{fig:rule-face-extrusion}) on a face}
  % \label{fig:before-after-triangulation}
  \label{fig:before-after-extrusion}
\end{figure}

In a few words, a \emph{rule} $r : L \longrightarrow R$ and a \emph{match} $m : L \to G$ to a G-map $G$, describe the transformation $G \longrightarrow^{r,m} H$ from $G$ to $H$.
$m$ allows to replace a sub-graph of $G$ described by the left-hand side of the rule $L$ with another one described the right-hand side $R$, on order to produce $H$.

Informally, in the triangulation rule shown in Fig.~\ref{fig:rule-triangulation}, the left side is made of only one node $n_{0}$ labeled with the face $\orbit{0, 1}$ type: this way, it can match any face.
For the match $m : n0 \mapsto a_{0}$ from $L$ to $G$ (Fig.~\ref{fig:house-triangulation-left-match}), the whole face $G\orbit{0, 1}(a_{0})$ is matched, highlighted in orange in the figure. $0$ (resp. $1$) links are represented by black (resp. red) segments.
On the right side, the node $n_{0}$ has label $\orbit{0, \_}$.
This means that when applying this rule, $0$ links of nodes $n_{0}$ have been preserved, while $1$ links have been deleted.
Hence, in $H$ (Fig.~\ref{fig:house-triangulation-right-match}), the $0$-links of the matched orange face have been preserved while $1$-links have been deleted.
In other words, the edges of the orange face are disconnected in $H$.
The new node $n_{2}$ of $R$ creates, in $H$, darts copied from the matched face.
This is why there are eight blue darts ($a_2$, $b_2$ \dots $h_2$) created from orange darts ($a_0$, $b_0$ \dots $h_0$).
Because $n_{2}$ is labeled $\orbit{1, 2}$, the  orange $0$ (resp. $1$) links on the left side of the rule are relabeled to $1$ (resp. $2$) links in the blue copy.
Therefore, the rule creates a dual vertex to the matched face.
Finally, the node $n_{1}$ creates a green copy with eight darts ($a_1$, $b_1$ \dots $h_1$), deletes the left-side $0$ links and relabels the left-side $1$ links to $2$ links.

All the highlighted links (Fig.~\ref{fig:before-after-house-triangulation}) are referred to as \emph{implicit links} in the rule nodes.
Conversely, the links connecting the nodes together in $R$ are referred as \emph{explicit links}.
For example, the explicit $0$ links between $n_1$ and $n_2$ link one-to-one the green and blue darts.
Therefore, the rule creates the four new edges $\orbit{0, 2}$ in the triangulation of the matched square face.
In the same way, the explicit $1$ links between $n_0$ and $n_1$ link one-to-one the orange and green darts.

The $n_0$ node (Fig.~\ref{fig:rule-triangulation}) is a \emph{preserved node} because it belongs to both the left and right side of the rule.
Nodes $n_1$ and $n_2$ are \emph{created nodes} because they belong only to the right side of the rule.
This rule does not comprise any \emph{deleted node} because none of its nodes belong only to the left side.

The \emph{orbit} notion is extended to patterns of rules.
For example, in the right-hand side of the triangulation rule (Fig.~\ref{fig:rule-triangulation}), the $\orbit{0, 2}$-orbit (an edge) incident to node $n_{0}$ contains the single node $n_0$, and the $\orbit{0, 2}$-orbit incident to $n_1$ contain the two nodes $n_1$ and $n_2$ and the explicit $\alpha_0$-link which connects them.
Additionally, an $\orbit{o}$-orbit is said to be \emph{complete} if each node in the orbit matches one link per label of $\orbit{o}$ either explicitly or implicitly.
For example, the $\orbit{0, 2}$-orbit incident to node $n_{0}$ in the right-hand side of the triangulation rule (Fig. 3) is not complete, because node $n_0$ has no $\alpha_2$-link, either implicitly or explicitly.
Conversely, the $\orbit{0, 2}$-orbit incident to $n_1$ is complete, because its two nodes $n_0$ and $n_1$ are incident to the explicit $\alpha_0$-link and both have an implicit $\alpha_2$-link.
Note that an $\orbit{o}$-orbit in a graph is entirely matched by a rule pattern if and only if the corresponding $\orbit{o}$-orbit in the pattern is complete.
For example, the $\orbit{0, 2}$-orbit incident to $a_0$ in the graph $H$ (Fig.~\ref{fig:house-triangulation-right-match}) is partially matched by the triangulation rule (only darts $a_0$ and $b_0$ are matched and not their $\alpha_2$ neighbors), because the $\orbit{0, 2}$-orbit incident to $n_0$ in the rule is not complete.
Conversely, the $\orbit{0, 2}$-orbit incident to $a_1$ in graph $H$ is entirely matched by the triangulation rule (both darts $a_1$ and $h_1$ are matched by node $n_1$ and darts $a_2$ and $h_2$ are matched by  node $n_2$), because the $\orbit{0, 2}$-orbit incident to $n_1$ in the rule is complete.

Jerboa's rules provide syntactic properties guaranteeing G-maps consistency preservation.

\subsection{Orbit tracking}

\begin{figure}
    \centering
    \includegraphics[width=\columnwidth]{rule-face-triangulation-3D-boxed.pdf}
    \caption{Face triangulation rule}
    \label{fig:rule-triangulation}
\end{figure}

\begin{figure}
\begin{subfigure}{.45\columnwidth}
    \centering
    \includegraphics[width=.8\textwidth]{cube-article-view.pdf}
    \caption{Graph H}
    \label{fig:demo-cube}
\end{subfigure}
\hfill
\begin{subfigure}{.45\columnwidth}
    \centering
    \includegraphics[width=.8\textwidth]{triangulated-cube-article-view.pdf}
    \caption{Graph I}
    \label{fig:demo-cube-triangulated}
\end{subfigure}
\caption{Application of the triangulation rule (Fig.~\ref{fig:rule-triangulation}) on the front left face}
\label{fig:concept-triangulation}
\end{figure}

Rules contains the necessary information characterizing the topological changes
affecting an orbit throughout an application. %
Thus, the tracking of an orbit is made automatic without any addition other than
the rules' syntactic analysis. %

For example, the application of the face triangulation rule
(Fig.~\ref{fig:rule-triangulation}) on the front left face of graph $H$
(Fig.~\ref{fig:demo-cube}), produces graph $I$
(Fig.~\ref{fig:demo-cube-triangulated}). %
The syntactic analysis of the rule allows us to deduce that the vertex
$\orbit{1,2,3}$ orbit (blue) incident to the node $n_2$ is \emph{created}
because $n_2$ itself is a created node. %
Moreover, this created vertex comes from the matched face $\orbit{0,1,3}$ orbit
incident to the node $n_0$ in the left-hand side of the rule, because the
implicit links $1,2,3$ of $n_2$ are relabeled from the matched implicit links
$0,1,3$ of $n_0$. %
Consequently, the corresponding orbit (blue) in the graph $I$ is created with
the matched face as \emph{origin}. %

In another example, the face $\orbit{0,1,3}$ orbit incident to the preserved
node $n_0$ is the whole pattern in the left-hand side of the rule and is the
pattern minus $2$-links (green) in the right-hand side. %
That is, the face orbit is \emph{split} along $1$-links because the second
implicit $1$-link of $n_0$ in the left-hand side is relabeled outside of the
face type orbit $\orbit{0, 1, 3}$ in the right-hand side. %
Consequently, the matched face of graph $H$ is split along its vertices'
$1$-links and into four faces in graph $I$. %

In a last example, the vertex orbit incident to the preserved node $n_0$ is
\emph{modified}. %
That is, the vertex orbit incident to $n_0$ in the left-hand side matches
$1$-links and because in the right-hand side the orbit comprises the node $n_2$,
the orbit comprises $2$-links. %
Since these $2$-links are relabeled from $n_0$'s $1$-links and there is no
split. %
Consequently, the matched vertices of graph~$H$ are modified in graph~$I$. %

With such analysis, topological changes can be logged within bulletin boards
which are automatically computed without requiring any other
intervention.

%%%% Local Variables:
%%% mode: latex
%%% TeX-master: "../main"
%%% End:
%%% Local Variables:
%%% mode: latex
%%% TeX-master: "../main"
%%% End:
